\documentclass[11pt,a4paper]{article}
\usepackage[italian]{babel}
\usepackage[margin=2.5cm]{geometry}
\usepackage{parskip}
\usepackage{amsmath, amssymb, amsthm}
\usepackage{booktabs}
\usepackage{array}
\usepackage{xcolor}
\usepackage{needspace}
\usepackage{enumitem}
\usepackage{listings}
\usepackage{tikz}
\usepackage[most]{tcolorbox}
\usepackage[hidelinks]{hyperref}
\usetikzlibrary{decorations.pathreplacing, shapes.geometric, arrows.meta, positioning}

% Niente vedove/orfane: i paragrafi non lasciano righe isolate a fine/inizio pagina
\widowpenalty=10000
\clubpenalty=10000

% Elenchi più compatti
\setlist{itemsep=2pt, parsep=0pt, topsep=4pt, partopsep=0pt}

% --- Riquadri con bordo nero, senza numero (si spezzano tra due pagine) ---
% Uso: \begin{definizione} ... \end{definizione}
%      \begin{teorema}[Nome] ... \end{teorema}  (il nome è facoltativo)
\tcbset{nota/.style={enhanced jigsaw, breakable, colback=white,
  colframe=black, boxrule=0.5pt, sharp corners}}
\NewTColorBox{definizione}{}{nota,
  before upper={\textbf{Definizione.}\ }}
\NewTColorBox{teorema}{o}{nota,
  before upper={\textbf{Teorema\IfValueT{#1}{ (#1)}.}\ }}
\NewTColorBox{proposizione}{o}{nota,
  before upper={\textbf{Proposizione\IfValueT{#1}{ (#1)}.}\ }}
\NewTColorBox{proprieta}{o}{nota,
  before upper={\textbf{Proprietà\IfValueT{#1}{ (#1)}.}\ }}
\NewTColorBox{assioma}{o}{nota, boxrule=1pt,
  before upper={\textbf{Assioma\IfValueT{#1}{ (#1)}.}\ }}

% --- Ambienti semplici (senza riquadro) ---
% Il titolo sta in cima al blocco, anche se il contenuto inizia con un riquadro o
% con codice; e non resta mai da solo in fondo a una pagina.
% Uso: \begin{esempio}[Titolo facoltativo] ... \end{esempio}
\newcommand{\newrunin}[2]{%
  \NewDocumentEnvironment{#1}{o}{%
    \par\Needspace{4\baselineskip}\addvspace{\medskipamount}\noindent
    \textbf{#2}\IfValueT{##1}{ (##1)}\textbf{.}\ \ignorespaces}%
  {\par\addvspace{\medskipamount}}}
\newrunin{esempio}{Esempio}
\newrunin{esercizio}{Esercizio}
\newrunin{osservazione}{Osservazione}

% V e F nelle tavole di verità
\newcommand{\V}{\textsf{V}}
\newcommand{\F}{\textsf{F}}

% --- Codice C ---
% Uso: \begin{codice} ... \end{codice}      codice C numerato
%      \begin{comando} ... \end{comando}    comandi da terminale
%      \begin{risultato} ... \end{risultato} output di un programma
%      \lstinline|printf("ciao");|           codice dentro al testo
\lstdefinestyle{c}{
  language=C,
  basicstyle=\ttfamily\small,
  keywordstyle=\color{blue}\bfseries,
  commentstyle=\color{gray}\itshape,
  stringstyle=\color{red!70!black},
  numbers=left,
  numberstyle=\tiny\color{gray},
  numbersep=8pt,
  columns=flexible,
  keepspaces=true,
  breaklines=true,
  showstringspaces=false,
  tabsize=4,
  morekeywords={include, define},
  % lettere accentate nei commenti
  literate={à}{{\`a}}1 {è}{{\`e}}1 {é}{{\'e}}1 {ì}{{\`i}}1
           {ò}{{\`o}}1 {ù}{{\`u}}1 {È}{{\`E}}1 {À}{{\`A}}1
}
\lstset{style=c}
\newtcblisting{codice}{listing only, nota, left=6mm,
  listing options={style=c}}
\newtcblisting{comando}{listing only, nota,
  listing options={basicstyle=\ttfamily\small, numbers=none, language={},
    columns=flexible, keepspaces=true, breaklines=true}}
\newtcblisting{risultato}{listing only, enhanced jigsaw, breakable,
  colback=black!5, colframe=black, boxrule=0.5pt, sharp corners,
  listing options={basicstyle=\ttfamily\small, numbers=none, language={},
    columns=flexible, keepspaces=true, breaklines=true}}

% --- Diagrammi di flusso (simboli standard) ---
\tikzset{
  terminale/.style = {ellipse, draw, minimum width=2.5cm, minimum height=0.9cm},
  io/.style        = {trapezium, trapezium left angle=70, trapezium right angle=110,
                      draw, minimum width=2.5cm, minimum height=0.9cm},
  processo/.style  = {rectangle, draw, minimum width=2.5cm, minimum height=0.9cm},
  decisione/.style = {diamond, draw, aspect=2, minimum width=2.5cm},
  freccia/.style   = {thick, -{Stealth}}
}

\title{Il campo dei numeri complessi}
\author{Marco Cerbella}
\date{Lezione 10 -- 20 ottobre 2026}

\begin{document}
\maketitle

\section{Il campo dei numeri complessi}

Consideriamo $\mathbb{R}^2 = \mathbb{R}\times\mathbb{R} = \{(x, y) \mid x, y \in \mathbb{R}\}$ con le operazioni
\begin{align*}
  (x_1, y_1) + (x_2, y_2) &= (x_1 + x_2,\ y_1 + y_2),\\
  (x_1, y_1)\cdot(x_2, y_2) &= (x_1 x_2 - y_1 y_2,\ x_1 y_2 + x_2 y_1).
\end{align*}
$(\mathbb{R}^2, +, \cdot) = \mathbb{C}$ è l'\emph{insieme dei numeri complessi}; $z = (x, y)$.

\subsection{Proprietà}

\begin{enumerate}
  \item \emph{$(\mathbb{C}, +)$ è un gruppo abeliano.}
  \begin{itemize}
    \item $z_1 + z_2 = z_2 + z_1$ (commutativa);
    \item $z_1 + (z_2 + z_3) = (z_1 + z_2) + z_3$ (associativa);
    \item esiste $0 = (0, 0)$ elemento neutro: $(x, y) + (0, 0) = (0, 0) + (x, y) = (x, y)$;
    \item per ogni $z = (x, y)$ esiste l'inverso (opposto) $-z = (-x, -y)$: $(x, y) + (-x, -y) = (x - x, y - y) = (0, 0) = 0$.
  \end{itemize}
  \item \emph{Distributive:} $z_1\cdot(z_2 + z_3) = z_1 z_2 + z_1 z_3$ e $(z_1 + z_2)\cdot z_3 = z_1 z_3 + z_2 z_3$.
  \item \emph{Associativa del prodotto:} $z_1\cdot(z_2\cdot z_3) = (z_1\cdot z_2)\cdot z_3$. Quindi $(\mathbb{C}, +, \cdot)$ è un anello.
  \item \emph{Commutativa del prodotto:} $z_1 \cdot z_2 = z_2 \cdot z_1$. Quindi è un anello commutativo.
  \item \emph{Elemento neutro del prodotto:} $1 = (1, 0)$ è l'unità per $\cdot$:
  \[
    (1, 0)\cdot(x, y) = (1\cdot x - 0\cdot y,\ 1\cdot y + 0\cdot x) = (x, y).
  \]
  Quindi $(\mathbb{C}, +, \cdot)$ è un anello commutativo con unità.
  \item \emph{Inverso moltiplicativo:} per ogni $z = (x, y) \neq (0, 0)$ esiste l'inverso
  \[
    z^{-1} = \left( \frac{x}{x^2 + y^2},\ -\frac{y}{x^2 + y^2} \right).
  \]
  Infatti
  \[
    (x, y)\cdot\left( \frac{x}{x^2+y^2},\ \frac{-y}{x^2+y^2} \right)
    = \left( \frac{x^2}{x^2+y^2} - \frac{y(-y)}{x^2+y^2},\ \frac{x(-y)}{x^2+y^2} + \frac{y\,x}{x^2+y^2} \right) = (1, 0) = 1 .
  \]
\end{enumerate}
Dunque $(\mathbb{C}, +, \cdot)$ è un \emph{campo}.

\subsection{Scrittura $a + ib$}

Per semplicità scriviamo $(1, 0) = 1$ e $(0, 1) = i$ (nuovo simbolo, \emph{unità immaginaria}). Allora
\[
  (x, y) = (x, 0) + (0, y) = x\,(1, 0) + y\,(0, 1) = x + yi .
\]
I numeri $(x, 0)$ sono i numeri reali: $(x_1, 0) + (x_2, 0) = (x_1 + x_2, 0)$ e $(x_1, 0)\cdot(x_2, 0) = (x_1 x_2, 0)$; quindi $\mathbb{R} \subseteq \mathbb{C}$.

\begin{osservazione}
$i^2 = i\cdot i = (0, 1)\cdot(0, 1) = (-1, 0) = -(1, 0) = -1$. Anche $(-i)^2 = (-i)(-i) = (0, -1)(0, -1) = (-1, 0) = -1$.
\end{osservazione}

L'equazione $x^2 + 1 = 0$ con $x \in \mathbb{R}$ non ha soluzioni. L'equazione $z^2 + 1 = 0$ con $z \in \mathbb{C}$, cioè $z^2 = -1$, ha soluzioni $z_1 = i$, $z_2 = -i$.

\begin{esempio}[Equazione di secondo grado su $\mathbb{C}$]
\[
  z^2 - 2z + 10 = 0, \qquad
  z_{1,2} = \frac{2 \pm \sqrt{4 - 40}}{2} = \frac{2 \pm \sqrt{-36}}{2} = \frac{2 \pm \sqrt{36\cdot(-1)}}{2} = \frac{2 \pm 6\sqrt{-1}}{2} ,
\]
cioè $z_1 = \dfrac{2 + 6i}{2} = 1 + 3i$ e $z_2 = \dfrac{2 - 6i}{2} = 1 - 3i$.
\end{esempio}

Possiamo scrivere i numeri complessi in questo modo: $z = a + ib$, $w = c + id$ con $a, b, c, d \in \mathbb{R}$; $a = a(1, 0)$ è la \emph{parte reale}, $b = b(0, 1)$ la \emph{parte immaginaria}. Allora
\[
  z + w = (a + c) + i(b + d),
\]
\[
  z \cdot w = (a + ib)(c + id) = ac + iad + ibc + i^2 bd = ac - bd + (ad + bc)\,i ,
\]
dove si è usato $i^2 = -1$.

\section{Riepilogo di esempi di strutture}

\begin{itemize}
  \item $(\mathbb{N}\setminus\{0\}, +)$: semigruppo (non monoide).
  \item $(\mathbb{N}, +)$: monoide (non gruppo).
  \item $(\mathbb{Z}, +)$: gruppo abeliano.
  \item $(2\mathbb{Z}, +, \cdot)$: anello commutativo non unitario.
  \item $(\mathbb{Z}, +, \cdot)$: anello unitario commutativo (non campo).
  \item $(\mathbb{R}, +, \cdot)$, $(\mathbb{Q}, +, \cdot)$, $(\mathbb{C}, +, \cdot)$: campi; anche $(\mathbb{Z}_2, +, \cdot)$ è un campo.
\end{itemize}

\end{document}
